\documentclass[11pt]{article}
\usepackage[T1]{fontenc}
\usepackage[utf8]{inputenc}
\usepackage[margin=1in]{geometry}
\usepackage{parskip}
\usepackage[colorlinks=true,urlcolor=blue,linkcolor=blue]{hyperref}
\pagestyle{empty}

\begin{document}

\noindent\textbf{Andrew H. Bond}, Senior Member, IEEE\\
Department of Computer Engineering, San Jos\'e State University\\
San Jos\'e, CA 95192, USA\\
\href{mailto:andrew.bond@sjsu.edu}{andrew.bond@sjsu.edu} \quad ORCID
\href{https://orcid.org/0009-0003-2599-6158}{0009-0003-2599-6158}

\medskip
\noindent July 2026

\bigskip
\noindent To the Editor-in-Chief\\
\emph{IEEE Transactions on Information Theory}

\bigskip
\noindent Dear Editor-in-Chief,

\medskip
I am pleased to submit the enclosed manuscript, ``Observation Theory: Geometry,
Distortion, and Reliability as Properties of Observation,'' for consideration as a
regular paper in the \emph{IEEE Transactions on Information Theory}.

The paper develops a consumer-relative theory of lossy representation. From a
second-order expansion of a downstream task's distortion it isolates a single
controlling quantity---the reconstruction error read through the consumer's read
operator $P_{\mathcal{C}}=J^\top G J$---and shows that ordinary mean-squared-error
reconstruction is its isotropic corner $P_{\mathcal{C}}=I$, the corner in which
Shannon's Gaussian rate--distortion function and Gauss least squares reside. The
principal contributions are: (i) an identification of the read operator with the
pullback geometry of \emph{distinguishability} and the quotient of the representation
it induces; (ii) the \emph{identifiability} of the read operator from black-box
consumer queries, validated blind against a planted ground truth; (iii) a
consumer-relative rate--distortion function with a proof of achievability and converse
for a general source; and (iv) a coding theorem for the \emph{true} nonlinear consumer
divergence that reduces exactly to output-side coding, together with its
high-resolution form. The framework is positioned as the identifiable geometric middle
term between Shannon rate--distortion and the information bottleneck.

The work lies squarely within the Transactions' scope---Shannon theory,
rate--distortion and source coding, quantization, and information geometry, with
connections to learning and inference. It draws on and extends the indirect/remote
source-coding tradition (Dobrushin--Tsybakov, Wolf--Ziv, Witsenhausen), coding for
computing (K\"orner, Orlitsky--Roche), and the rate--distortion--perception and
information-bottleneck literatures, all of which are cited and discussed.

The material is original and previously unpublished; it has not been presented, in
whole or in part, at any conference, and it is not under review at any other journal.
It relates to, but is distinct from, two companion works by the author (on the
spectral/angular observer and on geometric key quantization) that are cited in the
manuscript; the present paper is a self-contained information-theoretic synthesis and
is not a verbatim extension of either.

In the interest of reproducible research, every empirical claim in the paper is
governed by a sealed, before-the-fact registration with a pre-declared acceptance bar.
The pre-registrations, experiment code, committed result artifacts, and a
continuous-integration workflow that re-runs the reproducible experiments and verifies
the paper's theoretical constructions are publicly available at the repository named in
the manuscript. Consistent with the Society's encouragement, I intend to post a
preprint to arXiv at the time of submission.

I confirm that I am the sole author; that the submission is prepared in the required
single-column format and is well within the 50-page review limit; and that I hold a
registered ORCID. I would be glad to suggest qualified reviewers or to provide any
additional material the editors may find helpful.

Thank you for your time and for considering this work.

\bigskip
\noindent Sincerely,

\medskip
\noindent Andrew H. Bond, Senior Member, IEEE\\
Department of Computer Engineering, San Jos\'e State University\\
\href{mailto:andrew.bond@sjsu.edu}{andrew.bond@sjsu.edu}

\end{document}
